\chapter{重要重联问题与全文核实的新现象}\label{ch:findings}
\section{读结论之前，先区分四种不同问题}
本章只把完成正文、结果、讨论与关键图表核读的文献作为结论证据。尚未取得全文或尚未核读的文章仍在文献库和待核索引中，但不由摘要替作者下结论。各节末的文献键可在末章找到模型、参数、页码、图号与具体限制。

第一种问题是磁通能否改变拓扑；第二种是已启动的局地层处理磁通有多快；第三种是全系统多久完成重联；第四种是释放的能量如何进入各物种、热运动和小尺度。一个模型在第二种问题上成功，不意味着同时解决了其余三种。尤其要区分峰值率、准稳态率和固定时间窗平均率；局地上游场与初始全局场归一化也会改变比较结果。

\section{电子惯性重联：电流子层比电子惯性长度更薄}
\citet{Ottaviani1993}研究无耗散、不可压的二维冷两场模型，使用大 $\Delta'd_e$ 撕裂序列。其新发现不是简单的“电子惯性能断线”，而是\textbf{不同物理量出现不同层宽}：早期非线性阶段，速度与磁通扰动仍跨约 $d_e$ 的尺度，电流却产生不断缩窄的子层。广义磁通 $F$ 的拓扑仍被输运保存，普通磁通已改变。

论文用近似保持形状的宏观对流胞构造解析解。若 $A(t)$ 为流体位移幅度，在 $d_e\ll A\ll L$ 区间，子层厚度近似为 $d_e\exp[-A/\widehat d_e]$，其中 $\widehat d_e$ 与 $d_e$ 同阶；无量纲幅度方程含 $A''\sim A+cA^4$，体现非线性加速。这里的“准爆发”是有限适用区间内的加速，不能解释成真实等离子体产生无限电流。

本地JET预印本采用 $1024\times64$ Fourier分量，$L_x/L_y=0.5$、$d_e/(2L_x)=0.04$，仅一个线性不稳定模。作者在归一化时间约125时因子层无法解析而停止计算，并指出需要FLR、压力梯度、三维及可能的微不稳定性。这篇文章没有完成最终饱和和热化问题。

\section{暖电子约化模型：宏观饱和与微细结构继续演化共存}
\citet{Schep1994}构造的二维约化系统保存两族广义通量组合；其X点与普通磁通X点可以分离。有限电子质量使电流本身不再在相应定常传播解中奇异，但不保证电流梯度有界。该文的全FLR离子响应明确是满足若干极限与守恒要求的构造，而不是从Vlasov严格推出的任意尺度闭合；完整非线性本征解也没有求出。

\citet{L049}展示了这一结构的非线性后果。对于有限 $\rho_{\rm sound}/d_e$，$G_+$ 与 $G_-$ 在磁岛中沿不同方向卷绕，逐渐变成细丝。总能量仍守恒，磁能被转到小尺度流动和电子压缩等通道，粗粒化后的磁岛却可近似饱和。\textbf{“岛宽不再长大”与“所有自由度都停止演化”因此并不等价。}

冷电子极限则出现不同情况：初期电流更像条状层，随后速度剪切激发二次Kelvin--Helmholtz型不稳定性，产生单极与偶极涡、岛形脉动以及流体动能与电子动能间的振荡交换。随着 $\rho_{\rm sound}/d_e$ 增加到约1，所展示算例在岛饱和前不再出现该二次不稳定性。原文图1--3可直接比较暖电子细丝化、冷电子破碎和涡量结构。

这些图像不是任意参数下的相图。作者明确指出：无耗散暖电子层不断变薄，最终会使流体近似失效。随后用漂移动理学的正则动量薄片说明更丰富的相混合，但二维沿电流方向的对称性仍是关键条件。

\section{Hall解耦与快率：GEM解决了什么，没解决什么}
\citet{Biskamp1997}在准中性、均匀密度、标量压力的双流体模型中发现，电子层随电子尺度变化，而某些宏观通量率对电子质量和电子黏性弱敏感。计算包含正则化耗散，实际伪谱计算还采用六阶扩散。电子Kelvin--Helmholtz湍流主要关联有限耗散率，并没有显著增加原本已快的重联率。

该文表II按Alfvén时间归一化后，$d_i$ 从0.4降到0.05，峰值率由约0.68降到0.34；作者明确没有确定 $d_i\to0$ 极限。固定密度也移除了某些压力动力学。因此不能把它改写为“任意系统尺寸、无耗散且普适0.1”。

\citet{Birn2001}的GEM比较使用有限磁岛种子，刻意跳过线性触发阶段。Hall、混合与PIC模型的非线性通量曲线较接近，峰值约0.24；然而粒子模型中的电子电流层和离子喷流宽度可比对应流体模型大约两倍。文章还指出过量耗散可降低重联率，并把结论限定在二维。

因此GEM提供的是\textbf{同一已启动算例的部分宏观量可跨模型复现}的证据。它并未证明自然触发、压力张量、电子层厚度、能量分配和三维不稳定性均可由一个Hall项替代。

\section{色散波与开放排气：有力的模型内关联，有限的普适性}
\citet{Rogers2001}把whistler和kinetic Alfvén波的二次色散，与电子流出随层厚变化及排气开口联系起来。其等温两流体扫描中，两类分支都允许时重联率约0.12；仅一个分支时约0.06；强导引场、冷等离子体且两类二次色散分支均被抑制时，电流层较长，速率显著下降。

控制参数包括投影Alfvén速度、声速和质量比，而不是单独的 $B_g/B_{\rm rec}$。将线性分类用于非线性重联时，相关重联磁场分量自身随时间和位置变化；作者用上游反向场作方便的参数替代，并承认这一点。这不是适用于任意电子压力闭合的必要条件定理。下文电子CGL机制说明，同样的开放排气还应检查压力项而非仅凭波模名称解释。

2015年的两篇强导引场研究进一步给出\emph{同类两流体框架内}的反例：在$d_e\gtrsim\rho_s$、没有快速色散波的参数区，仍可得到对耗散弱敏感的快重联，并以PIC作对照。它们还重新检查2001年的扫描路径，发现启动期上游场耗尽与初始片宽也参与决定测得速率；详见本章第\ref{sec:withoutfdw}节。因而本节的早期结果应读作特定扫描中的关联，不能继续写成快重联的必要条件。\citep{L216,L217}

\section{快慢双稳态、迟滞与突发转换}
\citet{Cassak2005}的重要证据是同一参数、不同启动历史下出现两种状态。在其F3D算例中，归一化电阻约0.015时可得约0.06的快支与约0.01的慢支；降低电阻后再恢复，系统可留在快支。原文图3、4给出两个方向不同的转换阈值，这比仅比较两个不同参数算例更有力。

同时应记录，电子惯性长度为0.2而网格为0.1，只有两格跨一个 $d_e$；2005年正文没有单独展示网格收敛。对DIII-D的温度估算低约一个数量级，MRX比较当时缺少片宽数据。因此模拟中的迟滞证据与实验的定量验证必须分别评价。

\citet{Cassak2010}给出鞍结分岔的简化动力学解释，并在关闭电子惯性的Hall计算中仍得到双稳态。但它保留电阻和四阶耗散，并非无耗散证明；描述片宽变化的方程还没有自洽决定局地长度和演化中的上游场。该文对电阻平衡Hall态稳定性的论断，与下一节的大系统结果存在需要区分参数和几何的争议。

\section{磁岛先改变内层，再触发Hall阶段：中间态可以往复切换}
\citet{Baalrud2011}分清了两个阈值。长薄片的整体宽度可以仍大于 $d_i$，但撕裂扰动的电阻内层 $d_\eta$ 已小于 $d_i$，Hall项便可改变增长率与本征结构。适用排序为 $d_\eta\ll d_i\ll\delta_{\rm SP}$，不是整体电流片已经成为Hall开口层。

在该排序、快速增长和长波近似下，理论得到
\begin{equation}
{\gamma_{\max}\over\Gamma_0}\simeq3.2(d_i/L)^{6/13}S^{7/13},\quad
k_{\max}L\simeq2.9(d_i/L)^{1/13}S^{11/26},\quad\Gamma_0=2\VA/L.
\end{equation}
全非线性计算提取的初期增长率，在 $S=10^6$、$d_i/L\simeq2\times10^{-4}$ 附近从弱依赖变为明显依赖Hall尺度，与预测转折接近；作者也明确说，有限参数区和增长率提取误差不足以精确验证 $6/13$ 指数。图8显示长片中的小磁岛已带四极面外场。

\citet{Huang2011}研究另一层次：磁岛把长片切成短片，使局地片先达到Hall尺度。$S=5\times10^5$ 的扫描中，$L/d_i=2500$ 可保持开放单X点，5000时却在开放层、延长层和多磁岛状态之间循环，速率约在0.005--0.03变化。10000的情况没有进入相同Hall阶段。\textbf{“进入Hall尺度”不保证最终永久停在单X点。}

这些比较还受厚度定义影响：原文某时刻四分之一峰值全宽约 $0.7d_i$，半高半宽却约 $0.17d_i$。用不同定义讨论“片宽是否超过 $d_i$”会造成假矛盾。该文还说明，自然系统的尺度参数可能远离其循环中间区，不应把循环当作所有耀斑或磁尾的必然模式。

\section{不对称重联：三个特征点分开，标度与绝对数值分开}
\citet{Cassak2008}检验反平行、无导引场的非对称Hall重联。混合Alfvén速度为
\begin{equation}
v_{\rm out}^2\sim\frac{B_1B_2(B_1+B_2)}{\mu_0(\rho_1B_2+\rho_2B_1)}.
\end{equation}
其九组算例的通量率和喷流随不对称参数的变化大体符合该标度；对称基准率实际为约0.0587。值得关注的结构是：磁X点、电子停滞点和离子停滞点可以不同位置；磁岛向弱场一侧鼓起，电流片向强场一侧移动。测流入速度应使用移动X线参考系。

该文另一种离子流出速度定义会改变绝对数值约60\%，但不明显改变参数标度。说明“标度符合”比“所有流速数值相同”弱。计算还采用 $\beta_{\min}=4$、质量比25、惯性项中的常密度近似和四阶扩散。作者把可靠不对称程度限制在约4以内，并指出更大不对称时电子层会失真；不能直接外推到强密度差、导引场存在的磁层顶。

\section{电子层长度会依赖耗散，归一化会改变速率的表述}
\citet{Sullivan2009}在等温Hall／惯性双流体岛合并中扫描超电阻。无电子惯性时，电子层长约随 $\eta_H^{1/3}$ 变化，厚度约随 $\eta_H^{4/9}$ 变化；因此不能把局地理论中尚未求出的层长默认为只由边界决定的常数。

更值得学习的是原文图9、10的对照：用局地电子层上游场和密度归一化，速率随超电阻变化；用初始全局量归一化，或直接看原始电场，则速率近乎不变，约在0.1--0.2范围。这不是数据互相矛盾，而是局地归一化量本身随耗散改变。作者把“处理初始岛磁通所需的整体时间”与局地层的归一化几何区分开。

该研究\textbf{没有}重现动理学中长度约 $10d_i$ 的长电子片：其无惯性超电阻扫描最长仍不到一个 $d_i$。惯性扫描所得长度可拟合 $m_e^{3/8}$，但一数量级质量比区间不足以区分 $3/8$ 与 $1/3$。真实质量比下约 $15d_e$ 的长度和约一个 $d_e$ 的厚度均为外推。最低超电阻算例的厚度接近网格，原文也标明不可靠；不能仅凭论文题名中的“延长”断言已解决长电子层问题。

\section{磁岛合并：Hall模型的快峰值可能高估整体合并}
\citet{Ng2015}用同一Fadeev岛合并构型比较MHD、Hall、十矩与PIC。局地Hall层可以很快，但Hall平均率对岛尺度只有较弱依赖，而所比较PIC率约呈 $\lambda^{-0.8}$。Hall结果可高出PIC两到三倍。这里的平均窗是 $0$ 到 $1.5t_A$，不能与另一宽高比或不同平均窗的 $-0.5$ 指数直接混用。

十矩局地闭合通过调节离子松弛尺度可改善小系统压力支撑；大系统却产生PIC中未见的次级磁岛，并缺少相应弹回。\citet{Ng2017}改为非局地闭合后，片宽约 $1.20d_i$，接近PIC的 $1.15d_i$，而Hall约 $0.18d_i$；但平均率尺寸指数约 $-0.45$，仍不同于PIC约 $-0.8$，压力非旋转对称度也不完全匹配。

\citet{Ng2020}的温度梯度闭合进一步改善该尺寸依赖，并出现导引场从0增加到约1时岛合并峰值率先升、到2时再降的结果。不能用“导引场总会抑制重联”替代具体几何的计算。该研究的附录还说明，没有热流的某些稳态十矩停滞点会导致不合理的压力趋零；多演化六个压力分量并不会自动消除闭合难题。

\section{电子CGL的新机制：快率与四极场极性反转}
\citet{Cassak2015}逐项关闭Hall、电子惯性，或改用离子CGL、仅把电子各向异性放入动量方程，构成机制对照。特定电子CGL进入欧姆定律时，仍可得到约0.06--0.1的快率和开放排气；仅离子CGL或仅动量各向异性则较慢。面外四极扰动的极性与通常Hall四极场相反，是比“存在四极图样”更具体的判据。

不过模拟始终保留电阻0.005和四阶耗散，压力方程还省略电磁功加热；因此快几何机制不等于CGL压力在X点独自打破冻结。原文提到另一电子状态方程的慢结果只来自私人通信，PIC相符也未给出对照图。这些内容应作为进一步检验方向，不能与已展示的消融证据同权。

\section{非耗散双流体的数值证据与系统尺寸问题}
\citet{Andres2014}的无显式耗散伪谱比较提供了有价值的对照：提高分辨率时，理想MHD与无惯性Hall的伪通量变化趋小，保留电子惯性的模型则趋向非零变化。它比单独展示一张磁岛图更能支持惯性作用。

但总能误差小于 $10^{-8}$ 并不保证所有电子尺度结构准确。真实质量比、Hall参数0.1时，$d_e\simeq0.00233L_0$；$2\pi L_0$ 周期盒的2048网格间距约 $0.00307L_0$。全局通量相近仍不能称作充分解析一个 $d_e$ 宽的局地层。该研究的“惯性必要”限定在标量压力、不可压、无其他非理想项的模型中。

\citet{Andres2016}保持初始几何并扫描Hall参数0.07--0.16，共十组，测得电流层厚度约随 $d_e$ 变化，上游场和层长较弱变化，准稳态平均率约为 $0.11\lambda^{0.98}$。这是有限参数窗中的近线性趋势，不能简写成原归一化下 $\R=\lambda$；论文后一个关系另外把归一化长度和磁场改取局地值。

还须区分数值趋势与解析解释。该预印本的稳态推导有惯性旋度项符号不一致；严格对称、停滞的二维X点若无非对角压力和电阻，稳态对流惯性不能独自支撑有限 $E_z$，必须检查时间惯性或评估位置。报告保留可见的扫描结果，但不把该局地量纲推导当作已闭合的普遍稳态理论。

\section{三维Hall重联：湍流不一定比单X线更快}
\citet{Huang2024}使用压缩性Hall MHD、无电子惯性和标量绝热压力，在 $S=2\times10^5$、导引场与重联场同阶时比较二维和三维。其大算例达到 $2000\times1000\times2000$ 网格。二维两种Hall参数都形成单X点；三维中 $L/d_i=500$ 可接近二维态，1000则维持自生湍流。

新现象是\textbf{维数改变了最终状态选择}。并且单X点Hall态最快，三维Hall湍流虽快于相应电阻MHD，却慢于单X点Hall态；不能写成“湍流必然再加速”。文中三维率采用平均磁场代理量，需要与场线连接意义下的三维重联区别。亚离子谱区间还不够长，$k^{-11/3}$ 只是参考斜率，并非已证明的宽惯性区标度。

\section{十矩三维闭合：恢复不稳定性，同时产生新的误差}
\citet{Wang2015}的十矩与PIC比较不能概括成时间历程一致。两者在不同时间达到相同重联通量时，近X点一些结构相近；但速率随时间的曲线不同，某压力非对角分量峰值约相差四倍，下游压力结构也不一致。

\citet{Closure2026}在基于MMS非对称事件的三维设置中比较局地与梯度闭合。梯度闭合恢复了局地闭合所压制的低混合漂移和弯曲模，并产生次级磁岛、磁通绳和混合。计算采用质量比100、$1152\times576\times288$ 网格，梯度闭合开销约为局地方案三倍。

同时，梯度模型仍产生范围过大的各向异性与非旋转对称压力；部分峰值出现在与MMS／PIC不同的一侧；弯曲模在Vlasov中数个回旋周期即可发展，而该流体算例需数十周期。文章未提供网格与质量比联合扫描。它说明闭合会改变哪些模式能够出现，并未证明十矩已与动理学等价。

\section{嵌入宽电流片：漫长启动与局地快速阶段可以共存}
\citet{Shay2004}研究的初始电流片宽度随系统尺度变化，电子／离子扩散区则嵌在其中。七组二维周期算例中，四个较大算例在岛半宽约 $6.1$--$9.9d_i$ 后进入快速阶段；最大算例最后约3\%的模拟时长完成约90\%的磁通重联与75\%的岛宽增长。这个时间结构是本文的重要发现，不能概括成“从开始到结束都以0.1重联”。

作者实际测量局地上游场 $B_d(t)$：以离子面外电流峰值的25\%位置定义上游边缘，以电子和离子流出速度交点定义下游边缘；启动期另作修正。后期 $V_{\rm out}\propto B_d$，输入速度约为局地Alfvén速度的0.1。若以远处固定磁场归一化，同一时段的率会显著变化。研究没有扫描质量比，故其质量比无关的表述是联系先前工作后的预期，不是本组算例独立证明。

图8还给出一个容易被忽略的现象：小 $\Delta'$ 算例即使已呈开放X形，压力背压仍几乎抵消磁张力，离子喷流及重联仍很慢；较大 $\Delta'$ 的长片则在电子片变薄、Hall贡献增强后快速缩短。因此，单凭某一张X形磁力线图判断快重联并不可靠。默认 $d_e$ 只有两格，另一个高分辨算例还同时降低四阶耗散；这不是严格的独立网格收敛扫描。

\section{解析层模型：显式消去电阻，不等于已确定全局速率}
\citet{Malyshkin2008}从X点附近和电子层上游的方程展开得到连接Sweet--Parker与Hall极限的公式；强Hall极限为 $E/(B_m\VA)\sim d_i/L$。但 $L$ 是由上游磁场沿层方向的曲率定义的局地尺度，与 $B_m$ 一起作为输入。原文明确把求解它们留给全局问题。因此，该文解决了“给定局地几何后各项如何平衡”，没有从任意边界与初始电流片推出唯一的0.1。

\citet{Simakov2009}把电子区与离子区作有限区域平衡，以35组岛合并算例检验。其归一化率覆盖约1200倍，模拟／理论比在0.5--2.3之间；强Hall极限得到电子层长度 $w$ 控制的率。这里 $w$ 和局地上游场仍从模拟提取。重要的新结果是：\textbf{细长电子层可以对应对耗散系数不显含的率；开放电子X形也可能处于耗散控制的慢支。}该文的一般过渡公式是渐近匹配插值，不能称为没有近似的全局解析解。

\citet{Malyshkin2009}进一步把电子漂移惯性加入总动量，得到电子半厚约 $d_e$、电子层长度随有效电阻的倒数增长、率随有效电阻正比变化的稳态支。这里“collisionless”指电阻被解释为待动理学决定的有效量，X点仍满足 $E_z=\eta J_z$。它与 \citet{Andres2014} 的无显式耗散、非稳态惯性重联研究的是不同极限；两者并列并不构成逻辑矛盾。

\section{梯度闭合的推进：从直接热流检验到三维不稳定性}
\citet{L256}先从Vlasov分布函数计算热流矩，检查局地各向同性松弛闭合是否给出正确的散度形状、符号和幅度，再提出对压力偏离各向同性部分作Laplacian的闭合。这一步比仅比较重联率更有约束力：局地松弛会在实际热流很弱的上游区域产生错误的热流散度，改变 $k_0$ 不能同时修正所有分量的结构。

其梯度闭合在磁岛尺度 $\lambda/d_i=5,10,15,25$ 的算例中产生磁岛反弹，平均率约随 $\lambda^{-0.73}$ 下降；两套分辨率得到接近的平均与峰值率。但是GEM、大Harris片与岛合并分别用了 $k_{0s}d_s=3,1/3,1/2$。作者明确按动理学结果选取参数，故这些结果证明了模型的拟合与迁移潜力，还没有证明无需校准即可预测所有电流片。

\citet{E552}进一步使用统一的 $\chi=3,k_{0s}d_s=1$ 检验重联、岛合并和低杂漂移不稳定性（LHDI）。质量比100的特定静电LHDI算例中，Vlasov、梯度十矩、五矩和局地松弛十矩的增长率分别约为 $2.2,2.4,3.2,0.67\Omega_i$。这显示增加矩数本身不是充分条件：\textbf{五矩可以把不稳定性发展得过快，某些十矩闭合又会抑制过强。}最快增长模式会随模型和质量比变化，只检查预先指定的单一波数可能得出不同判断。

该文电磁LHDI支的Vlasov／流体增长率约为 $0.41/0.21\Omega_i$，都出现了超过初始重联场量级的磁扰动与强电流片弯曲，其高饱和原因仍不清楚。三维重联较二维更早启动，后期电流片加宽，出流区湍动强于电子扩散区；后期另一个不稳定性仅被称为“类似电流片剪切不稳定性”，尚未完成模态确诊。

一个尤其值得研究的新情况是\textbf{噪声决定不同过程的先后顺序}：密度扰动幅度由 $10^{-10}$ 增至 $10^{-2}$，独立测试的LHDI启动约从 $30\Omega_i^{-1}$ 提前到 $4\Omega_i^{-1}$。若不稳定性先长成，它能建立多个重联点；若种子重联先启动，又会改变密度梯度并减弱LHDI。因此，比较流体与PIC时，粒子噪声、初始种子与闭合导致的启动延迟都必须记账，不能把所有湍动差异都归因于“二维还是三维”。

\section{三维X线扩展：传播方向、启动延迟与有限磁能}
Shepherd和Cassak在电子携带全部初始电流的三维F3D算例中，观察到弱导引场主要沿电子运动方向扩展，较强导引场转为双向、接近导引场Alfvén速度。导引场为$2.5B_0$时，两侧测得约$3.0$和$2.9V_{\!A0}$。这不是无限精确的转变阈值：用电流阈值定义前沿有15\%--20\%速度不确定度，拟合还要截在斜向X线相互作用之前。弱导引场反方向的“零速度”只是该观测窗口内没有分辨。\citep{Shepherd2012}

Arencibia等把机制转成直接可检查的量：沿$z$变化的$E_x$是否在前方诱发正确极性的$B_y$。表I--III分别比较Hall电场梯度、对流电场梯度、磁张力与对应的磁场或速度变化，支持反平行时电子输运、强导引场时场线弯曲的解释。但研究还发现两个不能被简单传播标度省略的现象。\citep{L275}

第一，同样初始半厚$2d_i$，方形重联平面的X线只整体移动，而长度加倍的矩形域中X线持续伸长，速度约$-0.51V_{\!A0}$。作者提出方形域更早耗尽可用磁能；这是待能量与边界扫描检验的假说。第二，用压力种子替代磁场种子后，压力脉冲近磁声速度双向传播，但重联仍以约$-0.92V_{\!A0}$单向扩展。因此“最快传递的扰动”不自动等于“重联前沿”。此外，某切面$B_y$扰动约$t=60$到达，显著重联约$t=170$才建立。前沿判据及延迟必须和速度同时报告。

\subsection{追溯早期三维结果：有限长X线不是固定尺度定律}
Shay等的2003年模拟已发现：初始较厚电流片出现空间孤立的有限长X线，较薄电流片中的X线可逐渐连接，形成更广泛的磁通释放。所见约$10d_i$长度与磁尾BBF的量级比较是一种解释候选；粗电流方向网格、质量比25、缺少平衡法向磁场和微不稳定性，都限制了对亚暴的直接推断。将其与2021年长宽比试验联读，比单独引用“厚片形成孤立X线”更准确。\citep{Shay2003,L275}

\section{热离子改变细结构，未必明显改变总重联率}
Grasso等的三场陀螺流体模拟给出一个容易被只看磁通曲线漏掉的结果：热离子与热电子算例可有相近岛宽和发展速度，但静电势是否形成细丝、岛内密度是否变平、平行电场是否出现小尺度结构明显不同。热离子使静电势包含未经平滑的密度贡献，而陀螺平均势仍平滑；第一章已逐步推导这一关系。论文据此讨论流体电子束及可能增强的加速，但没有粒子分布函数，不能称为非热粒子谱的预测。\citep{E469}

另一个三维陀螺流体研究在小$\Delta'$参数下得到非线性爆发加速。最高扫描$\beta_e=1.6\times10^{-3}$时，以$L_\perp/c_s$归一的非线性峰值增长率为0.009，线性值0.0008。这里“增长率”来自扰动能量的对数导数，不能与$E/(B\VA)$直接比较。扫描同时改变盒子和初始片宽；X点分裂为两个Y点等拓扑变化在这篇短文中主要由文字报告，两幅结果图仅显示能量与增长率。三维是否必要、高速流是否是加速原因，作者仍留作后续问题。\citep{L184}

\section{递归片破裂：初级Hall很弱，次级Hall仍能显著增强}
Pucci等线性本征计算显示，Hall作用可以在$d_i<a$却$d_i>\delta_\eta$时改变撕裂，形成变宽的Hall中间层和变薄的内电阻层。拟合的增长指数与简单量纲匹配不同，并给出修正的临界长宽比；推导和适用范围已集中在第一章。\citep{L236}

Papini等的非线性计算进一步发现，初级磁岛之间自发形成半厚约初始$0.1$--$0.15$的次级片，片可先保持近压强平衡，再迅速破裂。按原片内层计$d_i/\delta_\eta=0.2$时，Hall几乎不改变初级增长，却使次级增长增加约17\%；比值1.4时次级增长接近加倍。关键新情况是：一个在宏观初级阶段可忽略的项，经由递归尺度缩小成为重要项。\citep{L262}

同时必须保留三条未决限制。第一，文中$\gamma_{\rm rec}$是对数磁通增长或磁扰动指数拟合，并非通常的$E/(B\VA)$；次级$2.3$与初级$0.48$不能与“普适0.1”直接比较。第二，换成各次级片自己的阿尔芬时间后，增长仍为一的量级；实测Hall次级片长宽比却\emph{不吻合}前述线性拟合标度，原文明确承认参数扫描不足。第三，强Hall下次级片不再可见，可能是更快破裂或X点几何改变，也可能是分辨率不足，不能写成已确证的“Hall必然抑制磁岛”。

\section{Hall模型也能使磁岛弹回，但结果受电子层处理控制}
Bard与Dorelli表明，不含离子动理学的Hall计算也可能在岛合并中出现弹回。在中等网格上，开放的双尺度层与弹回能够同时存在；继续细化后某些小系统不再弹回，而最大系统仍弹回并出现次级岛。这补充了早期“Hall合并过快”的结论：弹回并非动理学独占现象，但不同模型是否给出相同的平均完成时间仍需比较。\citep{L251,Ng2015,Ng2017}

该文以网格代替显式电阻，因此其层形转变并不是固定物理参数下发现了两个连续解。要把“数值耗散减少后的变化”提升为“物理双稳态”，还需要同物理参数、不同历史的持续分支检验；要建立普适合并标度，则需要控制电子层尺度和平均时间窗。

\section{电子独有重联：同样的能谱可隐藏不同的粒子参与}
Califano等的混合Vlasov计算以动理学离子和惯性等温电子比较两种湍流注入尺度。注入更靠近离子尺度的一组形成短电流片和电子独有喷流；另一组出现较长片以及离子、电子共同喷流。两组能谱却相似。这说明“谱看起来一致”不能单独验证流体模型的重联类型，必须同时检查离子响应、局地拓扑和片寿命。\citep{L267}

该文是\emph{动理学对照}而非纯Hall模拟。自动识别的数千电流峰并非全在重联，人工核对的重联层样本较小；电子惯性长度仅约1.36格，谱滤波与广义磁通的数值改变仍是限制。结果支持注入尺度可能选择电子独有重联的机制，不足以给出普适长度阈值。等温电子模型也不能用$\vct J\cdot\vct E'$直接预测电子升温。

\section{来源分群的新闭合方向：改善形状与解决热流并非同一结论}
2026年McGrae-Menge等预印本把片两侧电子分别演化为十矩流体。每群自身三阶矩较小时，群间相对漂移和温度差仍能产生很大的\emph{总}热流；第一章式\eqref{eq:mixture_heat}给出精确分解。与单电子十矩相比，两群模型改善了电子层密度空洞和压力堆积，重建的热流符号与空间结构更接近PIC，但幅度仍偏高。\citep{F635}

其证据边界尤其重要：离子固定；参数精度图是对各PIC算例在时间内选取\emph{最佳}重建指标，不是两流体模拟的全时段平均精度；不同图使用独立色标。高$\beta$时人工标签梯度可引发错误膨胀，晚期同群内部又出现多束。论文提出的动态分裂、三维、湍流及离子耦合尚未验证。所谓156倍压缩是表示自由度估算，不是测得的运行加速。这是一条具体且值得检验的新方向，尚不能写成一般热流闭合问题已经解决。

\section{密度空穴是额外检验：总率相近仍可能有显著结构误差}
Shay等2001年的跨模型研究不仅比较总率，还明确报告分离线密度耗尽：相同上游密度0.2下，Hall流体最低约0.03、混合模型约0.1、PIC约0.15。离子动理学喷流约宽两倍而峰速约减半，总输出通量因而仍接近。原文将空穴联系到Hall电流造成的磁场堆积和排斥离子的电场；有限轨道效应减弱空穴。\citep{Shay2001}

因此，“Hall模型算对了重联率”不能单独用来认证卫星沿轨道密度峰谷。该流体计算还为防止负密度加入二阶扩散，电子惯性用常$d_e$反演，使电子层比当地低密度PIC约窄两倍。其无质量电子计算靠六阶耗散断线，不是理想Hall模型自行改变磁拓扑。

该文关于小均匀电阻不能平衡电场的论证，使用了电子层宽长比近固定的几何关系。后续电阻Hall层理论放松几何约束，得到不同分支。阅读这些结论时，应先对照层长怎样决定，再判断是否互相矛盾。\citep{Shay2001,Malyshkin2008,Simakov2009}

\section{强场被带入后的自发启动：保持电阻也能转快}
Cassak等2006年从宽正弦片出发，发现慢重联把越来越强的上游场输送到内嵌SP层。电阻保持不变时，层仍逐渐变薄；达到当地$d_i$附近后离子电子解耦，出现开放排气与快速流入。慢阶段符合$\eta/\delta$，快阶段采用$0.17V_{\!A,\rm up}$拟合。这给出“先积累、再迅速释放”的一种具体模型过程。\citep{Cassak2006}

展示的低电阻运行早期由另一电阻运行重启，启动前瞬变已消失；因此准确说法是后期转变不需要在触发时继续改变电阻。太阳参数下的积累及转换时间只是量级相容，既未直接解析太阳的电子层，也未纳入高$S$磁岛、热反馈及三维竞争，不能据此唯一判断耀斑启动机制。

\section{强导引场磁绳合并：中间态阻塞，几何改变排岛}
MAST参数的Hall计算展示三类行为：厚于声回旋尺度的稳定碰撞片；中间尺度上迅速产生中心磁岛；薄层形成开放X点及高峰值率。平面算例保留180度对称，中心岛没有优先排出方向，长大后阻塞合并；环形轴对称几何打破这一对称，岛被排出并完成合并。新现象是外部几何可以决定同类局地岛是否终止全局重联。\citep{Stanier2013}

最低超电阻的峰值率较高，却有较慢的早期磁通处理，且开放分离线后来可再闭合、产生偏心岛对。超电阻控制断线同时加热局地电子，改变$\rho_s$；离子黏性也被人为增强。将这一组结果简化为“超电阻越小，重联一直越快”会丢掉文章最重要的非单调和时序信息。

\section{递归破裂后的三种状态：开放X点仍可持续排出磁岛}
Shi等在固定电阻、相同随机种子的七组Harris片计算中，观察到随$d_i$增强发生的变化：纯MHD反复形成更短电流片；弱Hall在开放X点继续产生并排出次级岛；较强Hall则形成不再明显生岛的开放层。排气开口角约由10度增至28度。这个中间态不是“Hall态与长SP片反复跳转”，而是开放结构与持续生岛共存。\citep{L263}

磁通斜率从约0.032升至0.058，但MHD晚期递归破裂后的0.034已经超过最弱Hall算例。较强Hall层长逐渐不变时，其半长仅剩约3--4个横向网格；最小MHD片也达到网格尺度。因此，观测到的状态差异比“已证实普适0.06上限”更可靠。它与Papini等的次级增长增强并不矛盾：线性增长、非线性开口和后期是否继续生岛，是三个阶段不同的问题。\citep{L262,L263}

\subsection{“Hall电流耗散”不是Hall项直接加热}
该文发现面内电流的电阻耗散$\eta(J_x^2+J_y^2)$占总$\eta J^2$的比例可从约40\%升到55\%，沿场电流的比例约45\%。这两种划分并不相互补足，因为面内／面外与平行／垂直是不同坐标分解。Hall项自身依然满足$\J\cdot(\J\times\B)=0$。变化的是Hall重排后的电流结构及其\emph{电阻}损失，不能称Hall项把一半磁能直接变成热。

原方程压强绝热演化，没有把$\eta J^2$回收到内能，因此这些量也不是完整的分物种加热预测。中心线磁谱从浅段加陡段过渡到近单一约$k^{-2.22}$，但固定线谱、有限拟合区间和不同平均时间窗还不足以证明普适湍流级联律。\citep{L263}

\section{没有快速色散波的快重联：反例具体推翻了什么}\label{sec:withoutfdw}
\citet{L216}在强导引场两场模型中分别研究$d_e>\rho_s$和$\rho_s>d_e$。前者没有通常用来解释快重联的快速色散分支。超电阻从$10^{-11}$增到$10^{-8}$时，两流体算例的归一化率仍在约0.025--0.029；对应单流体率则近于$\eta_H^{0.252}$。这是同一类耗散扫描中的实测差别，不是从X形图像直接判定快慢。

其关键新发现是\textbf{微观耗散层的上游场和长度会自调节}。当耗散层比$h=\max(d_e,\rho_s)$薄，上游$B_\delta$随厚度一起减小；某些区间内层长度也随厚度缩短。因此，对固定$B_\delta$和固定内层长度推导出的耗散依赖，不一定是实际算例采取的路径。第一章式\eqref{eq:two_layer_rate}把这一观察重建成两层匹配。它解释给定$B_h,w_v$时的弱耗散依赖，却没有自洽选出这两个外部量。

长文\citet{L217}进一步比较流体与PIC。在有快速色散波的一组参数中，两模型峰值率相差不到约10\%，但启动和饱和时刻不同；PIC早期出现的次级岛、喷流弯曲也没有在流体中逐一重现。在无快速色散波的\emph{小盒}算例中，没有次级岛仍能快速重联。因此，该组证据同时排除了“快速色散波必须存在”和“次级岛必须出现”这两种过强的必要性表述；并未排除它们在其他参数区影响速率。

\subsection{重做早期扫描：区分改变声尺度和改变电子尺度}
长文第VII节重做类似Rogers等的扫描：固定$d_e$而减小$\rho_s$时，速率下降可以出现；但固定$\rho_s$、在仍无快速色散波的区间增大$d_e$，速率也会上升。两组都改变了$h$相对于初始电流片宽的比例，因而改变启动等待时间及启动时还能供应多强的上游场。按当时测得的$B_h^2$归一化后，差别明显缩小。这里的结论是早期扫描不能单独隔离色散波的因果作用，并非“线性波对重联毫无作用”。\citep{Rogers2001,L217}

\subsection{尺寸扫描没有同时放大所有方向}
Harris系列把流出方向长度从基准的0.5倍变到8倍，基准$L_{x0}=100h$；入流方向盒宽与初始片厚保持不变。局地上游场取离中面约$4h$的位置，按入流穿越时间平均，所得流体与PIC率约0.09--0.17。长文还排除最后约四分之一时窗，以减少可用磁通耗尽的影响。这个结果支持所测试几何中的弱尺寸依赖，但不能替代全几何自相似放大，也不能证明任何边界下都为0.1。

学习时可先对照短文图1、3，再读长文图14--19：前者检查耗散层如何变化，后者检查尺寸、归一化和扫描路径。两篇包含相关算例，不能作为两次完全独立的统计验证。\citep{L216,L217}

\section{受迫重联：同一个“快”字包含不同驱动和结束方式}
\subsection{Taylor墙位移：快峰值后仍有漫长慢尾}
\citet{L057}研究最终停止移动的导体墙，给出线性和非线性、MHD和Hall的不同峰值标度。低电阻时，超电阻对X点电场的贡献会超过普通电阻；即使扫描中率对超电阻弱敏感，也不能把实际计算说成只含普通电阻的Hall模型。

任意导引场的四场扩展在$\beta_F=10^{-2}$至$10^2$范围、保持$d_\beta$相同时，得到相似的Hall突发阶段。保持的是$d_\beta=c_\beta d_i$，并非$d_i$固定。典型快峰值约为电阻参照的30倍，并有约8个Alfvén时间的延迟；\textbf{第一轮快爆发只把岛宽增长到最终值约四分之一}，随后仍有慢的Rutherford阶段。因此，“出现快速峰值”不等于“全部受迫磁通迅速完成重联”。\citep{L058}

两文将自身结果与近0.1的说法比较，但驱动、局地场和速率归一化不同。尤其后一文讨论中对上游磁场的线性标度与通常$B_uV_{Au}\propto B_u^2/\sqrt{n_u}$不同，不能只照搬其反驳措辞。更可靠的结论是：有限墙位移问题中的峰率和完成时间需要分别预测。

\subsection{体力持续压入：磁通堆积改变速率的分母}
\citet{L076}从较宽、撕裂稳定或边缘稳定的正弦片开始，用局域离子体力驱动入流。较弱驱动可建立平衡；较强驱动先积累上游磁场、压薄片，约到当地$0.7$--$0.8d_i$时出现快速阶段。其具体扫描中的临界驱动幅度在0.005--0.0075之间，这是选定盒子、驱动宽度和闭合的结果。

增加驱动宽度可提高\emph{绝对}电场，却因堆积的$B_u$更强而降低$B_uV_{Au}$归一化后的率。与自发薄片对照时，受迫算例的归一化出流和率可低约50\%。另一个值得单独记录的现象是：\textbf{磁岛宽度几乎饱和，岛内磁场却继续增加，故重联磁通仍增长}。用岛宽作为唯一通量代理会错判重联已经停止。

这篇2005年论文与2008年同题论文研究不同维数，不能合并成一个文件。2005年表I、正文算例交叉引用及部分图的驱动力数字有不一致；报告保留其主要可核实趋势，不为缺失输入自行补出精确对应。

\subsection{局域三维驱动：大部分磁通可以绕过重联层}
2008年研究把驱动同时沿第三方向局域化。二维中被驱动区域推来的上游磁通，超过70\%经过重联；相应三维例中约为10\%，其余可以沿第三方向偏转离开。三维局地电子电流还把磁结构向负$z$搬运，最大重联率不再位于驱动中心。持续显著重联的阶段也较短，典型比较约23对38个归一化时间。\citep{F717}

这里三维总磁通由场线追踪和面积积分求得，二维值则按单位第三维长度定义。作者将三维总率除以驱动全宽$2w_z=10d_i$作比较；这是一种有效长度归一化，不是一般三维重联率的唯一定义。典型对照的体力幅度也分别为二维0.15、三维0.2，并非仅改变维数的一对完全相同输入。结论应是\emph{局域第三维提供了磁通旁路，并改变局地层与整体驱动的联系}，而不是三维必然比二维快或慢一个固定倍数。

\section{陀螺诱导双重加速：先出现细尺度强流，再打开宏观排气}
\citet{F672}用强导引场四场陀螺流体，在大$\Delta'$、热离子和微观尺度远小于系统尺度时看到：线性增长之后先急剧加速，继而短暂减速或停滞，随后再次加速进入快重联。原文图6、8、9、11应连着看，分别追踪增长史、静电势、流场和排气开口；只看峰值会把两个过程合并。

第一轮加速与热离子有限回旋半径有关。在保持组合热尺度$\rho_\tau$相同、互换热离子／热电子贡献的对照中，热离子例出现更强的局地电势细结构，$E\times B$流可约高一个数量级；冷离子例缺少同样的第一轮加速。第一章已从准中性算子说明：热离子内区电势可直接跟随细密度结构，冷离子内区则经逆Laplacian平滑。这是与可测场量相连的机制对照。

第二轮加速随排气结构进一步开放，冷热离子均可出现。小尺度与宏观尺度分离越明显，两轮之间的减速越清楚。原文所画$\gamma_{\rm eff}=\dd\ln\delta\psi/\dd t$是\emph{对数磁通增长率}；它的尖峰不等于重联电场本身的尖峰，不能拿来与0.1直接比较。

表I将$d_e$从0.2减至0.025时，归一化峰值率由约0.283逐渐接近0.100；分母采用\emph{初始最大面内磁场}，不是随时间测得的局地上游场。作者也限定$d_e\Delta'\gtrsim1$：若固定$\Delta'$继续令$d_e\to0$，最终会进入小岛饱和区。因此该表不能证明严格理想极限仍有普适有限率。谱滤波、最细结构失去解析后的停止时刻和电子回旋效应缺失，同样限制了后期判断。

\subsection{哪些部分已经有动理学对照，哪些还没有}
\citet{E475}比较四场陀螺流体与EUTERPE陀螺动理学，先让两种代码各自对照其模型的线性本征解，再进行跨模型比较。这能把“离散是否解对自身方程”与“两个模型是否近似同一物理”分开。

在所扫线性波数中，增长率差别可达约20\%；减小平衡面内场后差别减弱，但没有自动降到零。非线性部分主要检验\emph{小}$\Delta'$：两模型都得到小岛饱和及其振荡，振荡频率很接近；饱和半宽差在较大驱动的$k_y=1.6$处约30\%，接近边缘稳定时明显减小。岛宽对温度相对不敏感，振荡频率却可随离子温度增加而近乎加倍。

论文文字称一个温比扫描的岛宽仅变约5\%，但图11显示的变化明显更大；据图估读，动理学约从0.040到0.049，流体约从0.046到0.053。这里只把它判为原文内部数量不一致，不把图上估读当作精确源数据。更关键的是，动理学非线性计算保留的模主要为$m=0,\pm1$，参数也在小岛区；\textbf{它没有检验前文大$\Delta'$的两次加速}。振荡可与粒子反弹频率比较，不代表有限矩模型恢复了完整俘获分布。

\section{加入Landau闭合：岛面积相近，增长时钟与温度结构已改变}
\citet{L244}构造含电子平行／垂直温度和离子垂直温度的六场模型，实际重联模拟则使用冷离子子系统。保留非局地电子热流后，线性增长下降；某些较高斜向模数的模式由不稳定变为稳定。闭合在谐振面局地为零，不妨碍它通过面外响应改变整个本征模。

非线性图像若用各自的$\gamma t$作横轴，岛形和最终面积可以相近，但实际时间已经不同。不能把这种重标度后的重合写成“Landau阻尼不改变重联速度”。温度、密度和电流的细结构也变平滑，所示电子平行温度峰值约降为无热流情况的一半。新发现是\textbf{宏观岛面积对闭合相对稳健，细结构和热响应却不稳健}。

其热流自由能损失可作为向未解析速度空间转移的代理，但模型没有复现所引用动理学研究中磁能释放与电子加热之间的明显时间延迟；作者提出闭合可能耗散过强。另一个量化限制是，约$0.002H_0$的数值能量误差已占本次释放磁能约20\%。故能量分配的精确比例没有像磁岛形态那样受到充分验证。

这些算例只含单一螺旋性，可以用$y^*=y+(k_z/k_y)z$约成二维；它们验证斜向结构与沿场闭合，却不能等同于多个斜向模相互作用的完整三维湍流。原文图4--7和能量段应一起读，避免只选最相似的岛图而忽略时间、闭合耗散与能量误差。

\section{百篇精读后的证据链：Hall效应改变了什么}
本轮39篇按模型争议、可复现细节和独立实验价值选择，而非按篇幅或容易下载来凑数。
以下侧重把新论文放回已有证据链；逐篇页码、参数和限制见第七章。

\subsection{早期闭合与移除Hall对照：不能只比较开关}
\citet{L029}的Hall/混合比较采用特定经验电子压力张量闭合，且其方程省略磁旋转项；这与今天完整十矩不是同一模型。
它提示离子非对角压力输运能改变层结构，不能从全局磁通相近推断离子应力无关。
\citet{L154}对无Hall快速重联的检查更具体：粗格Sweet--Parker层约2.5个网格时出现岛，约5格时可恢复较稳定层；去掉Hall后的率仍约0.015，含Hall约0.03，但仍有超电阻等耗散。
因此该对照揭示分辨率与耗散的作用，不能单独证明所有系统中Hall既必要又充分。

\citet{L035}的边界压缩算例中，混合模型因离子回旋效应出现更宽的响应，进入快速阶段的时间与Hall模型不同；局部嵌入Hall区可复现其特定总通量演化。
这解决的是\emph{该驱动条件下局地两流体区域是否足以改变整体响应}，未解决任意装置边界下的模型耦合。

\subsection{电阻从公式中消去，仍有未选定的长度}
\citet{L145}用电阻拖曳同时决定电子排出和耗散层厚，得到$E\sim B_mV_{\!A,m}d_i/L_e$；电阻消去后$L_e$和$B_m$仍是外部输入。
\citet{L137}从给定Harris外场及给定电场构造弱Hall局地解，也未选择全局率，且局部展开出现中心O点等几何问题。
\citet{L160}的有限质量比稳态分支又显示：当电子层厚受到$d_e$限制，$E\sim\eta B_m/(\mu_0d_e)$可显含电阻。
这些研究并不互相取消，而是采用不同闭合、长度约束和稳态假设；“快率与微观项无关”必须附带具体分支条件。

\citet{L089}从稳态不可压EMHD出发，用单位磁通的体积积分解释四极Hall场与电子流函数的关系。
该推导最有价值的细节是：理想EMHD内区的面外场主要由电子输运，生成需联系离子过渡区；预设X形面内场还在分离线产生对数奇异，必须靠惯性等修正。
它解释四极结构，未独立确定重联率和电子层长。\citet{L119}的实验只在远离X点和分离线的部分区域支持其沿射线近似恒定的场幅，近区明显偏离。

\subsection{线性Hall加速有条件，不能当成理想Hall直接断线}
\citet{L123}研究线绑边界的线性X点衰减；弱Hall拟合同时依赖$d_i$和电阻，因此将尺寸变化只归到$d_i^2$会漏掉归一化电阻的变化。
\citet{L122}提出三层变分近似，却没有完整$\Delta'$外匹配，且高斯试探电流导致某个$1/k^2$积分需额外正则化说明。其弱Hall修正还存在可直接代数复核的系数问题。
故本报告保留其启发，不将其形式上有限的无电阻增长解释成平滑理想Hall中磁通已被打破。

更系统的Hall撕裂分析\citep{L250}区分可压缩、常磁通与非常磁通分支；某一分支增长不显含$S$，仍需有限电阻完成内层匹配。
\citet{L238}在受迫稳定Taylor平衡中找到尤其值得注意的中间区：早期已出现Hall动力学，后期却能回到MHD式受限演化，尽管片厚小于$d_i$。
Hall还可能提高单模增长却增大最快波长，使可容纳的磁岛更少；“更快增长”与“更多磁岛”不是同一句话。

\section{三维并不只有“更快”：减速、倾角与传播记忆}
\subsection{电子片展宽与Hall场减弱也属于新现象}
\citet{L020}的三维惯性双流体计算中，重联建立的梯度激发静电涨落/低混合漂移类活动，电子层从约$d_e$展宽到$2.5d_e$，率比对应二维低约25\%，Hall场幅降低约30\%--40\%。
其平均惯性关联更像动量输运或异常黏性，不能统一称作异常电阻；密度梯度的放松还与磁场输运有关。
该算例质量比约278，$d_e$约3格，盒第三方向仅$1.25d_i$，因此是有条件的减速证据，不是三维必然减速定律。

\citet{L204}则发现Hall驱动的场向电流空间图样可不同于MHD的射流偏转机制；即使电流片较厚，系统沿电流方向的尺度与片厚之比仍会控制哪类图样主导。
模型没有地球偶极场、电离层反馈或电子加速闭合，不能直接声称重现完整亚暴电流楔。

\subsection{斜模更快与二维全局最快可以同时成立}
\citet{L224}在强导场force-free平衡、电子惯性和较弱超电阻下，固定波数比较发现小倾角撕裂比平行模快；导场比约10时部分倾角可快约两倍。提高超电阻又使平行模重新占优。
\citet{L268}使用电阻、不可压Harris平衡，完整扫描$(k_x,k_z)$后发现：某些固定较大$k_x$切片上斜模略快，但若盒能容纳较长波，全局最大增长仍在$k_z=0$。
两者的平衡、电子非理想项、扫描维度都不同，不能抽掉这些条件后写成互相证伪。
后者的复频率还表示扰动沿片传播，线性相速度不等于重联前沿扩展速度。

\citet{L166}三维不对称算例测得X线倾角约$35.9^\circ\pm8.5^\circ$，同时涵盖最大出流准则约$34.3^\circ$和最大率准则约$31.7^\circ$。
因此这是对角度量级的支持，精度不足以在两理论中作唯一选择；局地电流判据、有限盒和电子尺度分辨率也限制精确判别。

\subsection{重联沿非均匀电流片扩展会保留上游历史}
\citet{L294}得到薄片向厚片扩展时
$v_s\sim\VA d_iw_1/w^2$，厚片向薄片时约为$\VA d_i/w$。
前一种含起始薄片厚度$w_1$，因此前沿速度并非只由当前局地片厚决定；薄处生成的重联磁通进入厚处后被稀释，是这一“记忆”的来源。
这是比均匀片双向/单向传播更具体的新预测，适合与局域启动实验及非均匀磁尾比较。

论文表中的个别差值符号和磁尾应用时间需要独立检查。
按文内$w_1=0.89$或$3.5d_i$、$w_2=8.9d_i$、过渡指数2、距离$133d_i$、$d_i=720$ km、$\VA=1400$ km/s积分，相应传播时间约10.0--23.5分钟；不直接沿用正文2.5分钟的下限。
这是对该组输入的复算，不是对真实亚暴时间作新拟合；文后正式勘误涉及的参考文献/图例问题也不能自动视为已修正这些数值。

\section{增长加速、磁岛、涡旋和压强抑制的竞争}
\subsection{爆发律的更新伴随计算分辨率的更新}
\citet{L191}以能量变分纠正冷惯性模型的非线性爆发指数，\citet{L214}在暖双场模型中强调局域X几何及$\dot\epsilon\propto\epsilon^{7/4}$。
\citet{F697}进一步比较离子/电子温度与弱电阻：固定岛幅时某些速率和能量释放曲线对微观尺度较弱敏感，但从种子开始的总时间仍依赖线性内层尺度。
更关键的是，其冷模型会出次级岛并间歇减速；作者明确指出早期横向模数少和数值黏性漏掉了磁岛。
因此新结果是“局域爆发与磁岛间歇相竞争”，而不是一条适用于所有无碰撞重联的幂律。

\subsection{EMHD岛内部还会发生另一种不稳定性}
\citet{E459}在冷惯性EMHD中发现，早期有序的电流片级联被岛内射流的Kelvin--Helmholtz类涡旋破坏；电子压缩性可延迟或抑制这一过程。
它考察的主要是磁岛内部电流/涡量重新组织，不应只按“撕裂产生更多磁岛”概括。
模拟最高另做到$2048\times512$；晚期能量/Casimir误差可达5\%--10\%，不能由衰减曲线直接推断真实电子加热。

\citet{L202}则在惯性/电阻EMHD线性问题中强调纯增长本征模与驻波whistler不同，并构造有限中心电流的解。
其与MRX的增长时间和宽度只是尺度估计：计算约1.2 cm宽、实验比较约3.3 cm，且从假定10\%初扰动估总时间；未直接模拟稳态MRX排出和全欧姆预算。

\citet{L230}的强导场X点塌缩结果来自PIC，应放在动理学对照类。
其中电子涡旋可伴随次级X点，部分区域由对流惯性支撑电场；三维出现倾斜涡管。
但其三维图没有相应完整欧姆预算或全局率诊断，$Q$涡旋指标也不是重联判据，不能据图证明“电子压力项在三维全处消失”。

\subsection{磁岛触发Hall之后，重联还可能因压强结构被关停}
\citet{L159}在一组大$S$计算中看到Sweet--Parker、磁岛增强、Hall快速阶段先后出现，Hall阶段排出竞争磁岛。
后续文献已有开放X点持续排岛的情况，因此“Hall与磁岛互斥”只能限于该组边界、尺度与时段。
\citet{L174}则把压强梯度带入重联层，发现即使还有可释放磁能，反向抗磁漂移相对出流足够大时也可抑制重联。
该模型强导场、冷离子、绝热电子压力，MAST比较又依赖尺度估计；它提出了饱和/关停候选机制，尚非完整托卡马克锯齿预测。

\citet{E472}研究的是另一个问题：初始密度梯度使线性岛传播并降低增长；热离子时传播可朝离子抗磁方向。
其解析频率/增长率与数值有可见偏差，并未等价地验证\citet{L174}的非线性出流竞争条件。

\section{回旋平均、能谱与加热：全局相似不意味着模型等价}
\citet{E471}的四场陀螺流体显示，热离子改变导心流、X点尖角、电子密度分离线结构和高谐波，而部分全局磁通仍相近。
其$N_i$是导心密度，不是仪器测到的粒子密度；把两者直接比较会把极化响应误判成电子--离子非准中性。
\citet{E474}加入有限$\beta_e$和电子FLR后出现非线性减速间隔与第二轮超指数阶段，电子密度出现嵌套四极及更平缓梯度，而广义不变量仍继续细丝化。
此结果的方向随固定参数组改变，见式\eqref{eq:selected-beta}；部分增长率振荡还来自X点移动。

\citet{L264}比较Hall与混合动理学湍流时，磁场谱相似而速度及能量分配可不同；对比时段对应不同成熟阶段，不能仅由频率相近建立完整因果能量通路。
\citet{F644}的KH驱动Hall谱在离子尺度附近变陡，并在小电流片识别重联候选；但模型没有显式电子破冻结项，数值滤波就是必须审计的部分。
这些研究支持Hall模型捕获部分多尺度磁结构，尚未证明它精确预测粒子热化。

\citet{E512}提供了一个清楚的边界：电子平行相混合可把流体自由能送到高Hermite矩，再由弱碰撞热化。
这需要流体--动理学扩展，不能用纯Hamiltonian双流体的磁能下降代替。
\citet{E553}又从分布重建说明有限矩的内容边界：即使21矩恢复三角形主体，仍缺双峰与精细多轨道群。
\citet{L313}所讨论的电子正则磁场冻结则强调惯性拓扑变化不必耗散能量；它是理论框架讨论，没有本篇新的大规模重联模拟。
三维指数分离可缩短特定流程中的尺度形成时间，但三维并不自动保证混沌，常密度固定离子前提也不可随意推广。

\section{继续筛查发现的关键问题：相同速率不等于相同重联}
\label{sec:remaining-findings}
本节把此次15篇新增核读组织为问题链。它们包含直接Hall/双流体研究、混合模型、PIC、电阻MHD和实验，模型身份均保留；不能因用于Hall综述就把所有发现改称Hall模拟发现。

\subsection{全局重联是否存在系统尺寸效应}
\citet{F600}在固定宏观形状的二维Harris片中发现：尺寸增加30倍，PIC的全局平均率约降低2.6倍；Hall扫描增加100倍，率约降低3.4倍。小尺寸Hall结果近似$\ell^{-1/3}$，更大尺寸下降变缓，是否最终平台仍未确定。排气张角变窄伴随率下降；PIC持续产生磁岛，而Hall后期近单X点，Hall中的波列也没有在PIC中同样出现。

这对“跨尺度时一个固定率就足够”提出直接限制，但不构成对所有局地0.1机制的简单否定。初始片厚是否随盒子一起变、比较局地还是远上游场、按物理时间还是相同磁通阶段统计，会改变被回答的问题。第一章压力亏损模型针对局地准稳态；本篇针对整个演化的磁通处理，二者需要在同一诊断下重新比较。该文的平均区间排印矛盾及质量比25使其精确指数仍应审慎使用。

\subsection{导引场有时提高率，也会改变模型之间的差距}
磁岛合并研究\citep{L241}的结果比“导引场抑制重联”更丰富。小系统中从无导引场增加到$B_g=B_0$，PIC与混合模型峰值率分别提高约51\%和53\%，而Hall MHD略降。此时动理学模型中的离子更磁化，离子层变薄，堆积场和喷流增强，Hall与动理学模型的宏观率趋于接近。再增大导引场，则率重新缓慢下降。

无导引场时，PIC离子喷流比Hall结果弱约七倍，离子非旋转对称应力是重要阻力；Hall标量压力不能表达它。导引场较强时虽然宏观喷流相近，压力与应力的各项分解仍不相同。十矩用合适闭合长度可在小系统接近PIC，较大无导引场系统却产生PIC没有的次级岛，说明闭合的可移植性仍是问题。

全局岛合并还会反弹：若片内处理磁通不够快，两个岛先接近，再分开，之后再合并。无导引场PIC平均率随岛尺寸下降较陡；$B_g=B_0$时四种模型的平均率指数接近。大尺度无导引场PIC堆积区出现镜模相容的磁峰和密度凹陷；作者的线性色散检验支持该解释，但它是否进一步\emph{导致}重联减慢尚未证明。

\subsection{次级磁岛并不总带来更快的净能量释放}
三个证据应并列而非互相替代。第一，磁岛合并的强导引场PIC会不断生成次级岛，Hall却维持开放X形，二者平均率仍相近\citep{L241}。第二，PIC质量比扫描中低质量比容易生成岛，但岛并未使全局能量转换像电阻MHD算例那样显著提高；主电子层和移动次级X点的非理想预算也不同\citep{F606}。第三，欠分辨的MHD可以产生额外岛和假平台率\citep{F788}。

因此值得研究的新问题是：在什么驱动、寿命和尺度区间，岛的生成真正控制磁通通量？不能仅用“有岛”推断“岛介导的率”，也不能用一个无岛的已分辨算例否定其他初值下的撕裂。质量比趋于真实值、离子FLR、导引场、噪声和开放边界仍需联动检验。

\subsection{电子单独重联可以由大尺度重联和剪切流二次生成}
F655的2026年预印本\citep{F655}给出两条清楚的时序候选机制。其一，大尺度离子耦合重联先形成岛链，喷流使岛互相挤压，小于离子耦合尺度的新片再出现电子喷流。其二，湍流先生成电子速度剪切，随后出现与电子Kelvin--Helmholtz不稳定性相容的涡链，磁场被卷曲后形成小重联点。电子单独重联因此不必全靠直接在离子附近尺度注入能量。

这些图像有研究价值，但“太阳风中普遍发生”和“主导加热”尚超出本算例能证明的范围。系统仍为二维周期域，注入尺度约$50d_i$；电子尺度仅约一格。算法找到516个强电流结构，其中69\%薄于$d_i$且短于$5d_i$，并未逐个证实都是重联。部分$D_e$为负，所有片的平均量很小；不能把强局地正峰相加便当作净不可逆耗散。

结合解析速率研究\citep{F741}还得到另一个容易忽略的现象：无离子喷流并不意味着离子不反馈。周期小系统的离子反压仍可通过面内电场限制电子平台速度，使远上游离子归一率为$O(1)$；电子峰值、平台速度、局地电子归一率是三个不同量。

\subsection{Full-F、任意波长极化与多X点不同步}
GREENY及其后续论文\citep{E473,L314}显示，保留任意波长极化后，某些算例比长波近似发展更快并出现更细的电势和电流结构；全密度模型还保留了小振幅模型忽略的乘积项。初始电流由电子还是离子承担会明显改变过程，后续研究中离子载流启动通常慢约一个数量级。拉长盒子后可有多个X点在不同阶段同时活动，不能用一条单X点曲线代表整个域。

非正规线性算子提供暂态放大的可能，非线性阶段又伴随谱展宽；两者不能直接画等号。暖离子会改变态的形貌，但主岛数不随温比稳定单调增加。若只看摘要中的“新模型更快”或“非正规性导致爆发”，容易漏掉正文对范数、排序、初值和参数标签的限制。第一章给出了可以自行检查的能量范数和极化推导。

\subsection{Hall发电机既能协助，也能抵消磁松弛}
扩展MHD的RFP研究\citep{F798}不是孤立二维X点问题，而是三维多模撕裂和平均电流松弛。首次准单螺旋态崩解时，Hall与MHD发电机同向展平电流；后续多螺旋事件中，Hall项有时抵消部分MHD项。主模减弱、邻近模增强改变了相关电动势的符号；净松弛仍持续。

相应Lorentz力先驱动动量，Reynolds响应随后抵消一部分。后续事件中核心速度与电流的改变量同向，更接近MST集合观测；只看首次事件会得到相反判断。这个现象说明“Hall必然加速”过于笼统，也说明只验证总磁能释放相近不足以验证离子流和动量输运。

\subsection{多组分带来新的Hall图样和电场抵消}
冷离子MMS事件\citep{L227}表明，不同离子群可在同一分离线具有不同磁化程度；冷离子仍随电场漂移，其对流项可与Hall项同量级。负尘埃PIC\citep{L283}进一步出现内外极性相反的嵌套四极Hall场，外侧结构随着缓慢尘埃响应而建立。后者是三物种结果，不能列成标准电子--离子两流体已解释的现象；尘埃质量比、电荷和密度只取少数组合，也没有对应实验确认。

